\documentclass[11pt]{article}
\usepackage[margin=1in]{geometry}
\usepackage{amsmath,amssymb,amsthm}
\usepackage{graphicx}
\usepackage{microtype}
\usepackage[colorlinks=true,linkcolor=blue,citecolor=blue,urlcolor=blue]{hyperref}

\newtheorem{theorem}{Theorem}
\newtheorem{corollary}[theorem]{Corollary}
\newcommand{\Csim}{\mathcal C}
\newcommand{\SC}{\mathcal{SC}}
\newcommand{\T}{\mathcal T}
\newcommand{\Sg}{\mathcal S}

\title{A sharp quantitative tensor-splitting criterion without the zero-growth hypothesis}
\author{Automated mathematical discovery run \texttt{fa\_banach\_001}, lane 12}
\date{21 July 2026}

\begin{document}
\maketitle

\paragraph{Status.}
Strengthened partial result, likely valid, for human review.  It removes the
source theorem's zero-growth hypothesis when all but at most one factors have
growth-normalized similarity constant no larger than the tensor product's
similarity constant.  These factors may be nonnormal.  The general question
remains open.

\section{Source question}

Let $\T=(T(t))_{t\geq0}$ be a semigroup on a Hilbert space.  Write
$e_d\T=(e^{dt}T(t))_{t\geq0}$, let $\omega_0(\T)$ denote its exponential
growth bound, and let $\Csim(\T)$ be the optimal condition number of a
similarity taking $\T$ to a contraction semigroup.  Oliva-Maza and Tomilov
prove that if $\bigotimes_{k=1}^m\T_k\in\SC$ and its growth bound is zero,
then there are $d_k\in\mathbb R$ with $\sum_k d_k=0$ and
\[
  \max_k \Csim(e_{d_k}\T_k)
  \leq \Csim\!\left(\bigotimes_{k=1}^m\T_k\right).
\]
Immediately afterwards they ask whether the zero-growth assumption is
redundant \cite[Thm.~5.3 and p.~18]{OMT}.

\begin{center}
\includegraphics[width=.96\textwidth]{figures/open_problem_crop.png}
\end{center}

\section{Proof intuition}

The source paper's complete-boundedness proof contains more quantitative
information than its nonzero-growth splitting argument uses.  If one groups a
factor $\T_1$ against the tensor product of all remaining factors, then
normalizing the complementary tensor product to spectral radius one extracts
$\T_1$ with no loss in the similarity constant.  In general this produces an
endpoint shift incompatible with the endpoint shifts extracted for the other
factors.  The incompatibility disappears whenever the complementary factors
already have similarity constants at most the available product constant at
their growth-normalized endpoints: their endpoint shifts can simply be used,
and their sum supplies exactly the opposite shift for $\T_1$.  Contractivity
and normality are useful sufficient conditions, but they are not necessary.

\section{Strengthened partial theorem}

\begin{theorem}\label{thm:main}
Let $m\geq2$, and for $1\leq k\leq m$ let $\T_k$ be a nonzero semigroup on a
Hilbert space, strongly continuous on $(0,\infty)$.  Assume
\[
  \T:=\bigotimes_{k=1}^m\T_k\in\SC
\]
and put $C=\Csim(\T)$.  Suppose that the numbers
$\beta_k:=\omega_0(\T_k)$ are finite for $2\leq k\leq m$ and that
\begin{equation}\label{eq:endpoint-hyp}
  \Csim(e_{-\beta_k}\T_k)\leq C
  \qquad(2\leq k\leq m).
\end{equation}
Define
\[
 d_1:=\sum_{k=2}^m\beta_k,
 \qquad d_k:=-\beta_k\quad(2\leq k\leq m).
\]
Then $\sum_{k=1}^m d_k=0$ and
\[
  \max_{1\leq k\leq m}\Csim(e_{d_k}\T_k)\leq C.
\]
No assumption on $\omega_0(\T)$ is required.
\end{theorem}

\begin{proof}
Set
\[
  \Sg:=\bigotimes_{k=2}^m\T_k,
  \qquad \beta:=\omega_0(\Sg)=\sum_{k=2}^m\beta_k.
\]
The equality of growth bounds follows from the tensor-norm identity, as in
the source paper.  Since $\beta$ is finite, $\Sg$ is not quasinilpotent.

The endpoint estimate established in the alternative complete-boundedness
proof of the finite splitting theorem in \cite[Sec.~8]{OMT}, applied to the
two factors $\T_1$ and $\Sg$, gives
\begin{equation}\label{eq:endpoint}
  \Csim(e_{\beta}\T_1)
  \leq \Csim(\T_1\otimes\Sg)=C.
\end{equation}
Indeed, for every $t>0$ the operator
$e^{-\beta t}S(t)$ has spectral radius one.  The
Paulsen--Pearcy--Petrovi\'c endpoint estimate bounds the similarity constant
of $e^{\beta t}T_1(t)$ by $C$; Proposition 8.3 of the source then passes the
uniform pointwise estimate to the whole semigroup, yielding
\eqref{eq:endpoint}.

For $k\geq2$, hypothesis \eqref{eq:endpoint-hyp} says directly that
\[
  \Csim(e_{d_k}\T_k)=\Csim(e_{-\beta_k}\T_k)\leq C.
\]
Moreover $d_1=\beta$, so \eqref{eq:endpoint} supplies the remaining factor.
Finally,
\[
  \sum_{k=1}^m d_k
  =\sum_{k=2}^m\beta_k-\sum_{k=2}^m\beta_k=0.
\]
This proves the claim.
\end{proof}

\begin{corollary}\label{cor:normal}
The conclusion of Theorem~\ref{thm:main} holds if $\T_k$ is a normal
semigroup for every $2\leq k\leq m$ and each $\omega_0(\T_k)$ is finite.
In particular, for two factors the zero-growth hypothesis is unnecessary
whenever either factor is normal with finite growth bound.
\end{corollary}

\begin{proof}
For a normal operator, norm and spectral radius agree.  Therefore, for
$k\geq2$ and every $t>0$,
\[
 \left\|e^{-\beta_k t}T_k(t)\right\|
 =e^{-\beta_k t}r(T_k(t))=1.
\]
Thus each $e_{-\beta_k}\T_k$ is contractive and has similarity constant one,
so \eqref{eq:endpoint-hyp} holds.  Therefore
Theorem~\ref{thm:main} applies.
\end{proof}

For two factors, Theorem~\ref{thm:main} has the following useful form.  If
$\beta=\omega_0(\T_2)$ is finite and
\[
  \Csim(e_{-\beta}\T_2)\leq \Csim(\T_1\otimes\T_2),
\]
then the sharp conclusion holds with shifts $d_1=\beta$ and
$d_2=-\beta$.  This includes nonnormal factors whenever their endpoint
similarity can be implemented within the product's distortion budget.

\section{Finite-dimensional endpoint equality}

For a complex matrix $A$ and $C\geq1$, define
\[
 q_C(A):=\inf_{\substack{P>0\\\kappa(P)\leq C^2}}
 \frac12\lambda_{\max}\!\left(
 P^{-1/2}(A^*P+PA)P^{-1/2}\right),
\]
and write $s(A)=\max\{\operatorname{Re}\lambda:\lambda\in\sigma(A)\}$.

\begin{theorem}\label{thm:matrix}
Let $A\in M_n(\mathbb C)$, $B\in M_r(\mathbb C)$, $C\geq1$, and
$K=A\otimes I+I\otimes B$.  If $q_C(B)=s(B)$, then
\[
  q_C(K)\geq q_C(A)+q_C(B).
\]
The same conclusion holds with $A$ and $B$ interchanged.
\end{theorem}

\begin{proof}
Fix $\gamma>q_C(K)$.  After multiplying a feasible metric by a positive
scalar, there is $P>0$ such that
\[
 I\leq P\leq C^2I,
 \qquad K^*P+PK\leq2\gamma P.
\]
Choose a unit eigenvector $v$ of $B$ with eigenvalue $\lambda$ satisfying
$\operatorname{Re}\lambda=s(B)$, and compress $P$ to
\[
 P_v=(I\otimes v)^*P(I\otimes v)>0.
\]
Compression preserves the order bounds, so
$I\leq P_v\leq C^2I$.  Because $Bv=\lambda v$, compression of the Lyapunov
inequality gives
\[
 A^*P_v+P_vA+2\operatorname{Re}\lambda\,P_v
 \leq2\gamma P_v.
\]
Consequently $q_C(A)\leq\gamma-s(B)$.  Letting
$\gamma\downarrow q_C(K)$ and using $q_C(B)=s(B)$ proves the assertion.
\end{proof}

The equality hypothesis $q_C(B)=s(B)$ says precisely that the
growth-normalized matrix semigroup can be made contractive within distortion
$C$.  Theorem~\ref{thm:matrix} is therefore the infinitesimal matrix analogue
of Theorem~\ref{thm:main}.

\section{Verification, limitations, and the remaining obstruction}

The signs in the proof are forced: normalizing the complement uses
$e_{-\beta}\Sg$, while the endpoint extraction gives $e_{\beta}\T_1$.
The individual complementary shifts sum to $-\beta$, so the complete list
has zero sum.

The argument does not settle the general negative-growth case.  If neither
factor attains its spectral abscissa within the product's distortion budget,
the endpoint estimates supplied by the source may correspond to incompatible
zero-sum shifts.  The missing finite-dimensional implication is the
constrained-Lyapunov superadditivity inequality
\[
 q_C(A_1\otimes I+I\otimes A_2)
 \geq q_C(A_1)+q_C(A_2).
\]
This inequality is not proved here.  Convex SDP bisection described in the
accompanying attempt note tested more than 2,300 random real pairs in factor
dimensions two through four, in addition to structured Jordan families.  No
strictly active violation was found.  Numerical boundary artifacts were
discarded using the rigorous endpoint equality in
Theorem~\ref{thm:matrix}.  These computations are evidence only and play no
role in either theorem.

A bounded search on 21 July 2026 checked the original arXiv record, the June
2026 sequel on infinite tensor products, and keyword searches for the exact
redundancy question.  No later resolution of the general finite quantitative
question was located.  This packet claims only the stated partial theorem and
does not claim novelty beyond isolating the corollary and the remaining
metric obstruction.

\paragraph{Human-review recommendation.}
Verify the endpoint passage from Theorem 8.2 and Proposition 8.3 of the source
paper and the compression argument in Theorem~\ref{thm:matrix}; retain the
packet as a strengthened partial result.  Further work should focus on the
strictly active constrained-Lyapunov regime, where both factor endpoint
similarity constants exceed the available distortion budget.

\begin{thebibliography}{9}
\bibitem{OMT}
J. Oliva-Maza and Y. Tomilov,
\emph{On similarity to contraction semigroups and tensor products, I},
J. Funct. Anal. \textbf{291} (2026), no.~3, Paper 111492;
arXiv:2509.01005.
\end{thebibliography}

\end{document}
